%%%PAGE 167 rot=0 legibility=good summary=「粒子运动」提纲：电场+重力场、纯磁场（圆周/直线/螺旋）、组合场、叠加场、交变场；能量动量综合；找圆心四法；转筒类问题
%%%BLOCK id=167-01 ch=11 sec=带电粒子在复合场中的运动 kind=summary alt=09,04 cont=none y=0.08-0.55
\textbf{粒子运动}

① 电场$+$重力场\quad 等效重力场
\begin{itemize}
\item 类平抛/斜抛
\item 圆周
\item 直线 加减速
\end{itemize}

② 纯磁场
\begin{itemize}
\item $B\perp v$\quad 匀速圆周
\item $B\parallel v$\quad 匀速直线运动
\item $B$ 与 $v$ 夹角 $\ne 90^\circ$，$\ne 0^\circ$\quad 既有垂直分量又有水平分量\quad 等距螺旋线运动
\end{itemize}
%FIG p167_a 0.52 0.198 0.785 0.265
\notefig[0.4]{figures/p167_a.png}
螺旋半径
\[
R=\frac{mv_{\perp}}{qB}=\frac{mv\sin\theta}{qB}
\]
\[
T=\frac{2\pi m}{qB}\quad\text{与 $v$ 无关}
\]
\[
S=\frac{2\pi mv\cos\theta}{qB}=v_{\parallel}T
\]

③ 电磁组合拼接场\quad 速度关联、衔接，形状几何关系

④ 电磁复合叠加场\quad 分析\quad 无约束：直线、圆周、摆线。有约束：$N$ 方向会改变 $\to$ 具体分析

（右上方小注）直接分类讨论 $F$ 和 $qv_0B$ 的关系
% 上行为④行右端上方的小字注记；“的关系”四字写得更小，位于 qv_0B 下方
\par\smallskip
$\hookrightarrow$（自“无约束”一行引出）
\begin{itemize}
\item 直线必匀直
\item 圆周必匀圆
\item 周期由场定，与比荷无关
\end{itemize}
$T=\dfrac{2\pi m}{qB}=\dfrac{2\pi E}{gB}$\qquad $mg=qE$\qquad $F_{\text{洛}}=qvB=\dfrac{mv^2}{R}$
% 上述公式写在括号右侧，并有一斜线箭头指向“周期由场定，与比荷无关”

确定 $B$、$E$ 后，电性、运动方向也确定

⑤ 交变场
\begin{itemize}
\item $E$ 交变\quad $v$-$t$ 图
\item $B$ 交变\quad 按 $B$-$t$ 图画轨迹
\end{itemize}
由于 $B$、$E$ 有周期性，故粒子运动轨迹也有周期性、对称性、重复性
%FIG p167_b 0.355 0.466 0.52 0.545
\notefig[0.3]{figures/p167_b.png}
\[
\tan\theta=\frac{Eq}{mg}
\]
\[
\cos\theta=\frac{mg}{qvB}
\]
%%%END
%%%BLOCK id=167-02 ch=11 sec=带电粒子在复合场中的运动 kind=note alt=06,07 cont=none y=0.55-0.62
电磁场\ 能量\ 动量、力电综合，既有典型的力学过程，又有粒子运动的特点\\
\hspace*{2em}受力、运动、能量、动量
%%%END
%%%BLOCK id=167-03 ch=11 sec=有界磁场·找圆心 kind=method alt=none cont=none y=0.63-0.74
\textbf{找圆心}
\begin{itemize}
\item 出入速度作垂线 % 原稿此处先写“做”，其上方另写了“作”字（疑为改字）
\item 弦作中垂线 % “作”字处有涂改贴/白底
\item 偏转角作角平分线
\item 相切边界作垂线
\end{itemize}
%%%END
%%%BLOCK id=167-04 ch=11 sec=有界磁场·转筒类问题 kind=method alt=none cont=none y=0.76-0.81
转筒类问题\quad $\omega_{\text{粒子}}=\omega_{\text{筒}}$ 为不碰撞条件 % CHECK: 原稿“ω”字形近似“W”，据“转筒”语境读作角速度 ω
%%%END
%%%PAGEEND 167
